%==============================================================================
\section{Artificial Intelligence --- Trí tuệ nhân tạo và học máy}
%==============================================================================

\begin{dongco}[title={AI và ML: hai từ hay bị nhầm}]
``Artificial Intelligence'' và ``Machine Learning'' là hai cụm từ rất phổ biến trong công nghệ.
Chúng thường được dùng thay cho nhau nhưng \textbf{không} trùng nghĩa.
AI và ML có quan hệ với nhau nhưng khác về định nghĩa, ứng dụng và hàm ý.
Bài này so sánh sự khác biệt và mối liên hệ giữa hai khái niệm.
\end{dongco}

\subsection{Trí tuệ nhân tạo (AI) là gì?}

\begin{dinhnghia}[title={Định nghĩa AI}]
Trí tuệ nhân tạo là lĩnh vực rộng, phát triển các máy thông minh có thể thực hiện tác vụ thường cần trí tuệ con người: nhận thức, suy luận, học và ra quyết định.
Nói đơn giản: AI là khả năng của máy thực hiện tác vụ thường cần sự can thiệp hoặc trí tuệ của con người.
\end{dinhnghia}

\begin{phuongphap}[title={Hai loại AI}]
\begin{itemize}
  \item \textbf{AI hẹp (narrow / weak AI)}: thiết kế cho tác vụ cụ thể (nhận dạng giọng nói, nhận dạng ảnh, \ldots).
  \item \textbf{AI tổng quát (general / strong AI)}: có thể thực hiện mọi tác vụ trí tuệ như con người.
\end{itemize}
Hiện nay chỉ có AI hẹp đang được dùng; mục tiêu dài hạn là phát triển AI tổng quát áp dụng cho nhiều tác vụ.
\end{phuongphap}

\subsection{Các nhánh của AI}

\begin{ynghia}[title={AI như một ``giỏ'' các nhánh}]
AI gồm nhiều nhánh; các nhánh quan trọng gồm: Machine Learning (ML), Robotics, Expert Systems, Fuzzy Logic, Neural Networks, Computer Vision và Natural Language Processing (NLP).
\end{ynghia}

\begin{vidu}[title={Tổng quan quan hệ AI--ML}]
\tsimg{09_ai_neural_deep/ai_ml.jpg}
\end{vidu}

\subsubsection{Robotics}

\begin{ynghia}[title={Robotics}]
Robot chủ yếu được thiết kế để làm các tác vụ lặp lại, tẻ nhạt.
Robotics là nhánh AI nghiên cứu thiết kế, phát triển và điều khiển ứng dụng robot.
\end{ynghia}

\subsubsection{Computer Vision}

\begin{ynghia}[title={Computer Vision}]
Computer Vision giúp máy tính, robot và thiết bị số xử lý, hiểu ảnh/video kỹ thuật số và trích xuất thông tin quan trọng.
Với AI, Computer Vision xây dựng thuật toán trích xuất, phân tích và hiểu thông tin hữu ích từ ảnh số.
\end{ynghia}

\subsubsection{Expert Systems}

\begin{ynghia}[title={Expert Systems}]
Hệ chuyên gia là ứng dụng giải bài toán phức tạp trong một miền cụ thể, với trí tuệ, độ chính xác và chuyên môn gần như chuyên gia con người.
Giống chuyên gia, hệ này xuất sắc trong miền mà nó được huấn luyện.
\end{ynghia}

\subsubsection{Fuzzy Logic}

\begin{ynghia}[title={Fuzzy Logic}]
Máy tính thường nhận đầu vào số rõ ràng True/False; Fuzzy Logic là phương pháp suy luận giúp máy suy luận như con người trước khi quyết định.
Với Fuzzy Logic, máy có thể xét các khả năng trung gian giữa YES và NO, ví dụ ``Có thể đúng'', ``Có lẽ không'', \ldots
\end{ynghia}

\subsubsection{Neural Networks}

\begin{ynghia}[title={Neural Networks}]
Lấy cảm hứng từ mạng nơ-ron sinh học, Artificial Neural Networks (ANN) là tập các phần tử xử lý (node) kết nối dày, phản hồi động với đầu vào bên ngoài.
ANN dùng dữ liệu huấn luyện để cải thiện hiệu quả và độ chính xác (chi tiết ở bài Neural Networks).
\end{ynghia}

\subsubsection{Natural Language Processing (NLP)}

\begin{ynghia}[title={Natural Language Processing (NLP)}]
NLP cho phép hệ thông minh giao tiếp với con người bằng ngôn ngữ tự nhiên (ví dụ tiếng Anh).
Có thể ra lệnh cho robot bằng tiếng Anh đơn giản; NLP cũng xử lý văn bản và hiểu nghĩa đầy đủ.
Ứng dụng phổ biến: chatbot ảo, phân tích cảm xúc (sentiment analysis).
\end{ynghia}

\begin{vidu}[title={Ví dụ về AI}]
Trợ lý ảo, xe tự lái, nhận diện khuôn mặt, xử lý ngôn ngữ tự nhiên, hệ thống ra quyết định.
\end{vidu}

\subsection{AI và ML: tổng quan so sánh}

\begin{ynghia}[title={Bốn điểm khác biệt chính}]
\begin{enumerate}
  \item ML là \textbf{tập con} của AI; ML là kỹ thuật triển khai AI.
  \item ML tập trung thuật toán \textbf{học từ dữ liệu}; AI tập trung máy thông minh làm tác vụ cần trí tuệ con người.
  \item Thuật toán ML cải thiện độ chính xác theo thời gian nhờ dữ liệu; hệ AI có thể thích nghi tình huống/môi trường mới.
  \item ML thường \textbf{hẹp} hơn (chỉ học từ dữ liệu đã cho); AI hướng tới suy luận và quyết định trong bài toán phức tạp thực tế.
\end{enumerate}
\end{ynghia}

\subsection{Bảng: AI và Machine Learning}

\begin{vidu}[title={Khác biệt quan trọng}]
\begin{tabularx}{\linewidth}{>{\raggedright\arraybackslash}p{2.4cm}XX}
\textbf{Tiêu chí} & \textbf{Artificial Intelligence} & \textbf{Machine Learning} \\
\midrule
Định nghĩa &
Khả năng máy/hệ máy tính thực hiện tác vụ thường cần trí tuệ người (hiểu ngôn ngữ, nhận ảnh, ra quyết định). &
Một loại AI cho phép hệ học và cải thiện từ kinh nghiệm mà không lập trình tường minh; mô tả cách máy học và dùng kiến thức để cải thiện quyết định. \\
Khái niệm &
Xoay quanh thiết bị thông minh, thông minh. &
Xoay quanh việc cho máy học/quyết định và cải thiện kết quả. \\
Mục tiêu &
Mô phỏng trí tuệ con người để giải bài toán phức tạp. &
Học từ dữ liệu cho sẵn và cải thiện hiệu năng máy. \\
Bao gồm &
ANN, NLP, Fuzzy Logic, Robotics, Expert Systems, Computer Vision, ML, \ldots &
Supervised, unsupervised, reinforcement learning. \\
Phát triển &
Dẫn tới máy bắt chước hành vi con người. &
Giúp phát triển thuật toán tự học. \\
\end{tabularx}
\end{vidu}
